% Evans 12.3, printed pp.627--636; English PDF pp.644--653.
% Chinese start: printed p.444 / scan p.462.
\section{Semilinear Wave Equations}
\label{sec:12-3}
\subsection{A sign condition gives weak existence}
Assume $f$ is smooth and $zf(z)\ge0$. Then $F\ge0$.
The resulting energy controls $Du,u_t$, but it does not directly
control $\|u\|_2$ uniformly for all positive time.
\PDEAnnotation{能量不控制全时间质量}{这里 $F$ 可为零，能量只控制
$Du,u_t$。由 $u(t)=g+\int_0^t u_t$ 得到的是
$\|u(t)\|_2\le\|g\|_2+t\sqrt{2E(0)}$。
因此一般只能写每个有限时间条带上的 $L^\infty_tH^1_x$ 界。}
\begin{theorem}[Sign condition; source \S12.3.1, Theorem 1]
If $g\in H^1$, $h\in L^2$ and $F(g)\in L^1$, there is a global
distributional solution with $f(u)\in L^1_{\rm loc}$ and
\[
 \|Du(t)\|_2^2+\|u_t(t)\|_2^2+2\int F(u(t))\le2E(0).
\]
One can choose $u\in C([0,\infty);L^2)$ with
$\|u(t)\|_2\le\|g\|_2+t\sqrt{2E(0)}$, and $Du,u_t$ weakly
continuous into $L^2$. The energy inequality holds at every time
for these representatives. The initial energy traces are strong.
No uniqueness or energy equality is asserted in this general class.
\end{theorem}
\begin{expandedproof}[truncation, uniform integrability, and time slices]
Choose even smooth cutoffs $0\le\chi_j\le1$, equal to one on
$[-j,j]$ and zero outside $[-2j,2j]$. Put $f_j=\chi_jf$ and
$F_j(z)=\int_0^z f_j(a)da$. Then $f_j$ is globally Lipschitz,
$zf_j(z)\ge0$, $0\le F_j\le F$, and $F_j=F$ on $[-j,j]$.
The previous section gives an energy solution $u_j$.
Conservation for this Lipschitz problem follows by approximating the
data with smooth data in $H^1\times L^2$ and passing through the
linear energy estimate; $F_j(z)\le C_jz^2$ makes the potential
continuous in $L^1$ under $L^2$ convergence on bounded energy sets.
Consequently $E_j(t)=E_j(0)\le E(0)$ and
\[
 \|Du_j(t)\|_2+\|(u_j)_t(t)\|_2\le C_E,\qquad
 \|u_j(t)\|_2\le\|g\|_2+tC_E.
\]
Dominated convergence also gives $E_j(0)\to E(0)$.
Rellich compactness on bounded balls and the equicontinuity
$\|u_j(t)-u_j(s)\|_2\le C_E|t-s|$ give, by a diagonal
Arzela--Ascoli argument,
$u_j\to u$ in $C([0,T];L^2_{\rm loc})$ for every finite $T$,
and a.e. in space-time along a subsequence. The first derivatives
converge weak-star in $L^\infty_tL^2_x$.

Multiplying the approximate PDE by $u_j$ and integrating yields
\[
 \int_0^T\!\int u_j f_j(u_j)
 =\int_0^T\!\int((u_j)_t^2-|Du_j|^2)
      -\left[\int u_j(u_j)_t\right]_0^T\le C_{E,g,T}.
\]
All terms are justified first for smooth data and then by the same
energy approximation. The integrand on the left is nonnegative.
For a measurable subset $A$ of a compact space-time cylinder and $M>0$,
\[
 \int_A|f_j(u_j)|\le C_{E,g,T}/M
              +|A|\max_{|z|\le M}|f(z)|.
\]
The first term comes from $|u_j|>M$; the second uses $|f_j|\le|f|$.
This is uniform integrability. Also $f_j(u_j)\to f(u)$ a.e., because
at any point with finite limit $u$, eventually $|u_j|<j$.
Egoroff on the compact cylinder, followed by the preceding small-set
bound, proves $f_j(u_j)\to f(u)$ in $L^1$ there. This is a
space-time $L^1_{\rm loc}$ statement, not a claim of convergence
at every fixed time. The weak equation now passes to the limit.
For a smooth test reaching $t=0$ its first-order form is
\[
 \iint\{-u_t\varphi_t+Du\cdot D\varphi+f(u)\varphi\}
       =\int h\varphi(\cdot,0),\qquad u(0)=g.
\]

To obtain time representatives, for each compactly supported smooth
$\psi$ use
$\frac d{dt}\int(u_j)_t\psi=\int u_j\Delta\psi-f_j(u_j)\psi$.
Uniform integrability makes these scalar functions equicontinuous,
including at zero. Diagonal extraction for a dense collection of
$\psi$, and the uniform $L^2$ bounds, give weak $L^2$ convergence
of the velocities at every time and weak time continuity of $u_t$.
The gradients converge weakly at every time by the strong local
convergence of $u_j$. At a fixed time, extract a further subsequence
with a.e. spatial convergence on expanding balls. Fatou for $F_j(u_j)$
and weak lower semicontinuity of the two quadratic terms give
$E(t)\le E(0)$, independently of the chosen subsequence.
Thus this inequality holds for the chosen representatives at all times.
The identity $u(t)=g+\int_0^t u_t(s)ds$ gives the asserted global
$L^2$ continuity and growth bound. As $t\downarrow0$, lower
semicontinuity of the potential and weak convergence of $(Du,u_t)$,
together with $E(t)\le E(0)$, force convergence of their squared
norms; weak convergence plus convergence of Hilbert norms is strong.
This proves both strong initial energy traces. A diagonal extraction
in $T=1,2,\ldots$ completes global existence.
\end{expandedproof}

\subsection{Three-dimensional retarded potentials and continuation}
For smooth compactly supported $g,h$ let $v$ be their free wave.
In three dimensions Duhamel's formula is
\begin{equation}\label{eq:12-retarded}
 u(x,t)=v(x,t)-\frac1{4\pi}\int_{B(x,t)}
 \frac{f(u(y,t-|y-x|))}{|y-x|}\,dy.
\end{equation}
\begin{theorem}[Source \S12.3.2, Theorem 2]
If $f$ is smooth, $f(0)=0$, this problem has a unique local smooth
solution. If its maximal smooth lifetime $T_*<\infty$, then
\[
 \sup_{0\le t<T_*}\|u(t)\|_\infty=\infty.
\]
Equivalently, a bound for $u$ on a finite strip permits continuation.
We use this strip-bound criterion rather than infer a pointwise limit
from the mere failure of a uniform bound.
\end{theorem}
\begin{expandedproof}[the kernel norm and derivative persistence]
On $0\le t\le T\le1$, take the closed sup-norm ball
$\|u-v\|_\infty\le1$ of bounded continuous functions with initial
trace $g$. Its range is in $[-M,M]$ for a fixed $M$.
The right side of \eqref{eq:12-retarded} defines a map with
\[
 \|A[u]-A[w]\|_\infty\le\tfrac12 L_MT^2\|u-w\|_\infty,
 \qquad \|A[u]-v\|_\infty\le\tfrac12 T^2\max_{[-M,M]}|f|,
\]
because $(4\pi)^{-1}\int_{B(0,T)}|y|^{-1}dy=T^2/2$.
It is a contraction of the ball for small $T$.

The fixed point is smooth. For instance spatial difference quotients
satisfy a linear retarded equation with coefficient
$\int_0^1f'(\theta u(x+\varepsilon e_j,t)+(1-\theta)u(x,t))d\theta$.
It is uniformly bounded on a bounded range. The retarded operator
has the estimate
$\|R[G](t)\|_\infty\le\int_0^t(t-s)\|G(s)\|_\infty ds$.
Gronwall therefore bounds the first spatial derivatives. The time
derivative satisfies the differentiated wave equation with initial
data $(h,\Delta g-f(g))$ and is bounded in the same way. Successive
differentiation gives an equation with leading term $f'(u)D^\alpha u$
and forcing made from lower derivatives, already bounded by induction.
Difference quotients justify each differentiation and convergence;
the corresponding linear retarded equation has a unique continuous
solution. This bounds every derivative on a finite strip if $u$ is
bounded there. Finite propagation provides a common support ball.
Uniform bounds on one further derivative give limiting smooth data
at $T_*$; applying local existence to them extends the solution,
a contradiction. Alternatively higher-order energy limits provide
the same restart. Uniqueness follows either from contraction on
successive short strips or the difference-cone energy estimate.
\end{expandedproof}

\subsection{Subcritical defocusing powers}
Consider $f(u)=|u|^{p-1}u$, $1<p<5$, in $n=3$.
\PDEAnnotation{实数幂与光滑性}{当 $p$ 不是正奇整数，
$|z|^{p-1}z$ 在零点通常不是 $C^\infty$。
下面对所有 $1<p<5$ 声称有限能量有界 mild 解；
只有非线性本身光滑时，才由导数持续性得到光滑解。}
\begin{lemma}[Cone-trace bounds; expanded source \S12.3.3]
For a smooth defocusing solution and $t\le T_0<\infty$, put
$w(y)=u(y,t-|y-x|)$ on $B(x,t)$. Then
\begin{equation}\label{eq:12-cone-hardy}
 \|Dw\|_2+t^{-1}\|w\|_2+\|w/|y-x|\|_2+\|w\|_6
 \le C(E(0)^{1/2}+T_0^{1/2}\|g\|_\infty).
\end{equation}
Constants are independent of the apex and the cone radius.
\end{lemma}
\begin{expandedproof}[the nonzero boundary value in Hardy's inequality]
The chain rule gives $Dw=Du-\nu u_t$ on the mantle.
Thus \eqref{eq:12-cone-flux} gives $\int|Dw|^2\le2E(0)$.
Let $\overline w$ be the ball average. Integrating $|u_t|$ along
vertical lines in the cone gives
\[
 \int_{B(x,t)}|w|\le\int_{B(x,t)}|g|
                        +\iint_{K_{0,t}}|u_t|.
\]
Indeed each vertical segment from time zero to $t-|y-x|$ lies in
the cone. Since its space-time volume is $Ct^4$ and
$\iint_Ku_t^2\le2tE(0)$,
$|\overline w|\le C t^{-1/2}E(0)^{1/2}+\|g\|_\infty$.
Poincare's inequality yields
$t^{-2}\|w\|_2^2\le C\|Dw\|_2^2+Ct|\overline w|^2$.
The scaled ball inequalities are
\[
 \|w\|_6+\|w/|y-x|\|_2
 \le C(\|Dw\|_2+t^{-1}\|w\|_2).
\]
They follow by scaling the unit-ball Sobolev and Hardy inequalities.
For the latter, extend the unit-ball function by a bounded $H^1$
extension and cutoff to $\R^3$, then apply whole-space Hardy.
One cannot apply the zero-trace version directly to this cone trace.
These inequalities and the average bound prove the lemma.
\end{expandedproof}

\begin{theorem}[Subcritical existence; source \S12.3.3, Theorem 3]
For compactly supported smooth data and $1<p<5$, there is a unique
global finite-energy mild solution, bounded on each finite time
strip, with $u\in C_tH^1$, $u_t\in C_tL^2$ and conserved energy.
When the nonlinearity is smooth, the solution is smooth.
\end{theorem}
\begin{expandedproof}[the sublinear bootstrap and the power qualification]
First work in a smooth regime. Formula \eqref{eq:12-retarded} and
Cauchy--Schwarz give
\[
 |u(x,t)|\le\|v\|_\infty+C
 \|w/|y-x|\|_2\left(\int_{B(x,t)}|w|^{2(p-1)}dy\right)^{1/2}.
\]
For $1<p\le4$, the last exponent is at most six. Holder on the ball
and \eqref{eq:12-cone-hardy} give a uniform bound for $t\le T_0$;
this uses integrals, not an invalid norm triangle inequality when
$2(p-1)<1$. For $4<p<5$, put
$M(\tau)=\sup_{x,0\le t\le\tau}|u(x,t)|$. Then
\[
 \left(\int|w|^{2(p-1)}\right)^{1/2}
 \le M(\tau)^{p-4}\|w\|_6^3,
 \quad M(\tau)\le C_0+C_1M(\tau)^{p-4}.
\]
Since $0<p-4<1$, Young's inequality absorbs the last term into
$M/2$ plus a constant. Thus $M$ is bounded independently of the
current smooth lifespan below $T_0$. The continuation theorem
proves global smooth existence whenever $f$ is smooth.

For an arbitrary real $p$, approximate $f$ by smooth sign-preserving
nonlinearities, cut off at large amplitudes, converging in $C^1$
on every bounded interval and with
$|f_j(z)|\le C(|z|+|z|^p)$. For example smooth
$z(z^2+\varepsilon_j^2)^{(p-1)/2}$ first and then truncate far out.
Their nonnegative potentials and compact initial data have uniformly
bounded energy. The preceding argument gives a uniform strip bound;
the extra linear term is bounded by the same Hardy estimate times
the square root of the ball volume. On this bounded range the
nonlinearities have one uniform Lipschitz constant. Difference energy
estimates, using uniform convergence of $f_j$ on that range and the
common finite support ball, show Cauchy convergence in
$C_tH^1\times C_tL^2$. The retarded Volterra difference estimate
also gives Cauchy convergence in the supremum norm.
The retarded formula passes to the limit,
and its sup-norm contraction gives the same mild solution. The
potentials converge uniformly on the bounded range, so their integrals
converge and energy equality passes to the limit. The difference
estimate also proves uniqueness among bounded energy solutions on
each finite strip. Smoothness across zeros is not claimed for the
noninteger powers. This supplies the regularity qualification absent
from a literal all-power smoothness statement.
\end{expandedproof}
